\documentclass{article}

\usepackage[final]{neurips_2020}

\usepackage[utf8]{inputenc}
\usepackage[T1]{fontenc}
\usepackage{hyperref}
\usepackage{url}
\usepackage{booktabs}
\usepackage{amsmath}
\usepackage{amssymb}
\usepackage{amsfonts}
\usepackage{nicefrac}
\usepackage{microtype}
\usepackage{graphicx}
\usepackage{subcaption}
\usepackage{xcolor}

\graphicspath{{images/}}

\newcommand{\dataset}{\textsc{VlaExec}}
\newcommand{\datasettrain}{\textsc{VlaExec-train}}
\newcommand{\datasettest}{\textsc{VlaExec-test}}

\title{
  MemoryGrounding: Mid-Execution Grounding Validity Monitoring \\
  for Language-Guided Robot Manipulation \\
  \vspace{1em}
  \small{\normalfont Korea University COSE461 Final Project}
}

\author{
  Seongmin, Lee\; Sanghyuk, Ko \\
  Department of Computer Science \\
  Team 26 \\
  2020320055\;   2020320035
}

\begin{document}

\maketitle

\begin{abstract}
Vision--language models used for ambiguity resolution in language-guided
robot manipulation evaluate the scene only at task initiation. We show
that this static formulation fails when the scene changes during
execution: a representative initial-scene baseline
(AmbResVLM~\cite{ambresvlm}) achieves $14.3\%$ decision accuracy on
mid-execution scenarios. We introduce \textbf{MemoryGrounding}, a
two-stage framework that (i)~monitors video with a lightweight finite
state machine over Grounding DINO~\cite{gdino} (G-DINO) detections to
trigger grounding re-evaluation, and (ii)~issues structured
\textsc{Continue}/\textsc{Ask}/\textsc{Stop} decisions through a
Molmo-7B~\cite{molmo} evaluator conditioned on the stored initial
grounding $G_0$, the current detections, and the checkpoint image. On
\datasettest (70 trials over 7 scenario types covering
count change, displacement, and disappearance at the pre-pick and
pre-place checkpoints), the full system attains $100\%$ decision
accuracy with $0\%$ miss rate and $0\%$ false-alarm rate, together with
$90.0\%$ strict accuracy that additionally requires the clarifying
question to characterize the change type correctly. Controlled ablations
isolate three sources of gain: G-DINO supplies quantitative scene
perception, $G_0$ as temporal memory is required for displacement
detection ($+31.4$\,pp over a fine-tuned model without $G_0$), and
task-specific fine-tuning corrects force-matching bias where rule-based
pipelines fail on all target-removal trials. A counterfactual diagnostic
that varies $G_0$ while keeping the image and detections fixed yields a
memory reliance rate of $100\%$ ($41/41$), providing direct evidence
that the model relies on $G_0$ rather than count-based shortcuts.
\end{abstract}

\section{Introduction}
\label{sec:intro}

Language-guided robot manipulation has progressed rapidly with
vision--language models (VLMs) that ground natural-language instructions
in visual scenes~\cite{molmo,rt2,openvla}. A central challenge is
\emph{ambiguity}: when an instruction such as ``pick the cup and put it
in the red box'' does not uniquely identify the target or destination,
the robot must request clarification before acting. Existing methods
address this at task initiation---\emph{before} the robot
moves---by detecting whether the initial scene is ambiguous with respect
to the instruction~\cite{ambresvlm,knowno,clara,lap,sgcot}.

Robot manipulation is, however, a temporally extended process. Between
task acceptance and the irreversible pick or place action, the scene
can evolve: a second same-category object may appear, the original
target may be displaced, or the target may be removed entirely. A scene
that was unambiguous at task initiation can therefore become ambiguous
or invalid mid-execution, and existing methods provide no mechanism for
detecting or responding to these dynamic changes. We organize our
investigation around three research questions.

\textbf{RQ1.} \emph{Can VLM-based ambiguity detectors perceive
quantitative scene changes during execution?} Methods such as AmbResVLM
treat ambiguity as a binary judgment, ignoring how many target instances
are present and their spatial relation to the original plan. On our
mid-execution benchmark AmbResVLM attains only $14.3\%$ decision
accuracy, failing on all count-change and disappearance cases. This
motivates integrating Grounding DINO~\cite{gdino} (G-DINO) as a
supplementary recognition module that supplies count and coordinate
information at each checkpoint.

\textbf{RQ2.} \emph{Is prior state information necessary to detect
certain classes of grounding change?} When a target moves to a new
position the count remains $1$ but the identity may have changed; the
current frame alone provides no evidence of motion. A fine-tuned model
without access to the initial grounding $G_0$ (FT-A) attains $68.6\%$
accuracy and fails on all displacement scenarios (target moved $0/10$,
destination moved $0/10$). Storing
$G_0 = \{(l_\text{tgt}, \mathbf{x}_\text{tgt}),\
       (l_\text{dst}, \mathbf{x}_\text{dst})\}$
as a temporal reference enables comparison against the baseline and
yields $100\%$ accuracy (FT-B). On our benchmark $G_0$ thus constitutes
the minimal temporal memory sufficient for mid-execution grounding
validity monitoring.

\textbf{RQ3.} \emph{Does fine-tuning overcome force-matching bias under
fixed task queries?} Open-vocabulary detectors return the
highest-confidence match for a text query regardless of whether the
target is actually present. When the target has been removed, G-DINO
may match ``cup'' to a visually similar object, reporting $n=1$ at an
unexpected position. A rule-based system that trusts this detection
outputs \textsc{Ask} (ambiguous) instead of \textsc{Stop} (missing),
failing on all $10$ target-removal trials. Task-specific fine-tuning
enables the model to cross-validate detections against the image and
recover the correct \textsc{Stop} decision ($10/10$).

We address these questions with \textbf{MemoryGrounding}, a two-stage
monitor. \textbf{Stage~1} runs a lightweight finite-state machine (FSM)
over G-DINO detections on the video stream and triggers re-evaluation
when count shifts, displacement, or disappearance occur.
\textbf{Stage~2} is a Molmo-7B evaluator that, upon trigger, takes the
checkpoint image, the per-label G-DINO coordinates, and the stored
$G_0$, and outputs a structured \textsc{Continue}/\textsc{Ask}/\textsc{Stop}
decision with a clarifying question. For controlled comparison with
fixed-checkpoint baselines we evaluate Stage~2 in isolation at
pre-defined checkpoint images, consistent with the protocol of prior
work.

Our contributions are: (i)~an empirical analysis showing that existing
VLM-based ambiguity detection fails mid-execution due to two missing
capabilities---quantitative scene perception and temporal memory;
(ii)~an FSM-based vision trigger that fires only on semantically
significant mid-execution scene changes; (iii)~a $G_0$-conditioned
fine-tuned Molmo-7B evaluator whose ablation isolates the three sources
of gain; and (iv)~the \datasettest benchmark of $70$
samples across $7$ scenario types.

\section{Related Work}
\label{sec:related}

\paragraph{Pre-execution ambiguity resolution.}
AmbResVLM~\cite{ambresvlm} uses a two-step Molmo~\cite{molmo} protocol
to detect ambiguity and generate clarifying questions at task
initiation. KnowNo~\cite{knowno} casts ambiguity as conformal prediction
over actions with statistical task-completion guarantees but at the
action-selection level, without object-level grounding.
CLARA~\cite{clara}, LAP~\cite{lap}, and SG-CoT~\cite{sgcot} use language
models or scene-graph reasoning for pre-execution clarification.
INVIGORATE~\cite{invigorate} maintains a candidate graph and eliminates
candidates through dialogue but in a static interaction paradigm. None
of these methods detects grounding invalidation that occurs after motion
has begun.

\paragraph{Mid-execution monitoring.}
The Physical Agentic Loop~\cite{pal} monitors execution at transition
points using physical signals (gripper telemetry, joint torque) to
detect failures---its checkpoints are analogous to our $C_1$/$C_2$, but
its signal is \emph{post-action} and \emph{physical}: it asks ``did this
grasp physically succeed?''. We ask a complementary question, ``is the
original semantic grounding still valid?'', and do so \emph{before}
irreversible actions. DoReMi~\cite{doremi} uses LLM reasoning over robot
state to detect post-action failures; like PAL, it is reactive to
execution outcomes rather than semantic scene changes.

\paragraph{Trigger detection for scene change.}
Behavior-tree disambiguation~\cite{btdisambig} embeds disambiguation
into BT preconditions, triggering AmbResVLM when a same-category object
appears. Our FSM generalizes this beyond same-category addition to the
full taxonomy of mid-execution changes (addition, removal, displacement
at both target and destination) and routes the trigger into a
memory-conditioned evaluator.

\paragraph{Detector force-matching and visual verification.}
G-DINO~\cite{gdino} assigns each text query to the most query-like
region regardless of presence, an inherent force-matching behavior.
Molmo~\cite{molmo} is trained to point at objects by coordinate and
tends to return empty results when an object is genuinely absent; we
exploit this complementary signal via a fine-tuned Molmo-based evaluator
that cross-validates G-DINO detections.

\paragraph{Adaptation.}
Our evaluator is obtained by LoRA~\cite{lora} fine-tuning of Molmo-7B
for structured grounding-state classification, rather than action
generation as in RT-2~\cite{rt2} and OpenVLA~\cite{openvla}.

\paragraph{Positioning.}
Table~\ref{tab:related} summarizes the closest systems. In contrast to
prior work, MemoryGrounding combines FSM-based vision triggers,
force-match--aware ensemble evaluation, and $G_0$-conditioned decision
making within a single framework.

\begin{table}[h]
\centering
\small
\caption{Positioning against the closest related systems.
\textsc{Trig}: vision-based trigger.
\textsc{Eval}: structured grounding evaluator.
\textsc{$G_0$}: temporal reference memory.
\textsc{FM}: force-match correction.}
\label{tab:related}
\begin{tabular}{lcccc}
\toprule
System & Trig & Eval & $G_0$ & FM \\
\midrule
AmbResVLM~\cite{ambresvlm}      & N & init only & N & N \\
KnowNo~\cite{knowno}            & N & init only & N & N \\
PAL~\cite{pal}                  & physical & N & N & N \\
BT-Disambig.~\cite{btdisambig}  & same-cat. & limited & N & N \\
INVIGORATE~\cite{invigorate}    & N & static & partial & N \\
\textbf{MemoryGrounding (ours)} & \textbf{FSM} & \textbf{ensemble+FT} & \textbf{Y} & \textbf{Y} \\
\bottomrule
\end{tabular}
\end{table}

\section{Approach}
\label{sec:method}

\begin{figure}[t]
\centering
\includegraphics[width=\linewidth]{pipeline.png}
\caption{Overview of \textbf{MemoryGrounding}. \emph{Stage~1} (left):
the video stream is processed by G-DINO and a lightweight
trigger/checkpoint monitor that issues re-evaluation requests on count
shifts, displacement, or disappearance. \emph{Stage~2} (right): upon
trigger, a structured prompt assembling the task, the stored $G_0$
(initial target and destination labels and pixel coordinates), and the
current G-DINO detections is fed together with the checkpoint image to a
LoRA-fine-tuned Molmo-7B evaluator, which emits a state-and-question
JSON consumed by a deterministic decision policy.}
\label{fig:pipeline}
\end{figure}

\subsection{Problem formulation}

We consider a language-guided pick-and-place task with instruction
$\mathcal{L}$. At task initiation ($t_0$) the robot establishes an
\emph{initial grounding}
\[
G_0 = \{\text{target}: (l_\text{tgt}, \mathbf{x}_\text{tgt}),\
        \text{dest}:   (l_\text{dst}, \mathbf{x}_\text{dst})\},
\]
where $l$ is a text label and $\mathbf{x}\in\mathbb{R}^2$ a pixel
centroid. \textbf{Stage~1} monitors the video stream $\{I_t\}$ for time
points where scene change warrants re-evaluation. \textbf{Stage~2}
receives the triggered checkpoint image $I_{t_k}$, the stored $G_0$,
and the current G-DINO detections, classifies the grounding state, and
issues a decision
$d\in\{\textsc{Continue},\textsc{Ask},\textsc{Stop}\}$ accompanied by a
clarifying question (Table~\ref{tab:taxonomy}).

\begin{figure}[t]
\centering
\includegraphics[width=0.95\linewidth]{fixed_checkpoint_pipeline.png}
\caption{Fixed-checkpoint evaluation protocol. At $t_0$ the system
stores $G_0 = \{(l_\text{tgt},\mathbf{x}_\text{tgt}),
(l_\text{dst},\mathbf{x}_\text{dst})\}$. At $C_1$ (pre-pick) the
evaluator classifies the target state into \textsc{Clear},
\textsc{AmbiguousTarget} (which subsumes displaced targets), or
\textsc{InvalidTarget}; at $C_2$ (pre-place) the destination state is
classified analogously. Each state maps to a
\textsc{Continue}/\textsc{Ask}/\textsc{Stop} decision.}
\label{fig:fixed-checkpoint}
\end{figure}

\begin{table}[h]
\centering
\small
\caption{Grounding state taxonomy and decision policy.
\textsc{InvDest} maps to \textsc{Stop} when the destination is absent
with no visible substitute ($S_6$), and to \textsc{Ask} when an
alternative location may be negotiated.}
\label{tab:taxonomy}
\begin{tabular}{lll}
\toprule
State & Trigger condition & Decision \\
\midrule
\textsc{Clear} & objects at expected positions, unique & \textsc{Continue} \\
\textsc{AmbTarget} & multiple targets or displacement ($C_1$) & \textsc{Ask} \\
\textsc{AmbDest} & multiple destinations or displacement ($C_2$) & \textsc{Ask} \\
\textsc{InvTarget} & target absent ($C_1$) & \textsc{Stop} \\
\textsc{InvDest} & destination absent ($C_2$) & \textsc{Ask}/\textsc{Stop} \\
\bottomrule
\end{tabular}
\end{table}

\subsection{Stage 1: FSM trigger}

A lightweight FSM operates on per-frame G-DINO detections and triggers
re-evaluation on (i)~\emph{count increase}, $|D_t|>|D_\text{ref}|$;
(ii)~\emph{position shift}, $\|\mathbf{x}_t-\mathbf{x}_\text{ref}\|>\tau$;
or (iii)~\emph{disappearance}, $|D_t|=0$. The reference record
$r_\text{ref}$ is fixed between triggers and updated only after firing,
which prevents accumulated drift; a $k$-frame cooldown suppresses
repeated firing on the same event. On the video sequences in
\datasettrain, the FSM produces no false triggers on the
\textsc{Continue} scenario and successfully triggers on every
non-\textsc{Continue} scenario. We therefore treat Stage~1 as a solved
module and evaluate Stage~2 at pre-defined $C_1$/$C_2$ checkpoint
images---a protocol consistent with prior work and necessary for fair
comparison with AmbResVLM, which operates on static images.

\subsection{Stage 2: G-DINO, $G_0$, and fine-tuning}

\paragraph{G-DINO as supplementary recognition.}
AmbResVLM treats ambiguity as a binary judgment without tracking count
or position, which limits it to $14.3\%$ accuracy on our benchmark.
G-DINO returns, for each queried label, the centroids of all detected
instances, allowing discrimination of $n=0$ (absent), $n>1$ (multiple
candidates), and $n=1$ at a position consistent with $G_0$
(\textsc{Clear}) versus inconsistent (displaced). The displacement case
relies on $G_0$ as a reference.

\paragraph{$G_0$ as temporal reference.}
On our benchmark, $G_0$ together with the current G-DINO detections is
sufficient to disambiguate all seven scenario types via two binary
tests---count change and position shift. The reverse implication also
holds: for displacement scenarios ($S_4$, $S_7$), $n=1$ at both $t_0$
and the checkpoint, so the current frame alone is indistinguishable
from \textsc{Clear}. Only by comparing the current detection against
the stored $t_0$ position can the system infer motion.

\paragraph{Fine-tuned evaluator with force-match correction.}
Even with $G_0$ and G-DINO, a rule-based pipeline attains $72.9\%$
because G-DINO force-matches: when the target has been removed, it
returns a spurious detection on the most query-like region. A
rule-based system reads $n=1$ as ``present'' and outputs \textsc{Ask}
instead of \textsc{Stop}, failing on every target-removal trial. We
therefore train a Molmo-7B evaluator that consumes the checkpoint
image, the stored $G_0$ in text form, and the per-label G-DINO
coordinates, and produces a JSON object
$\{\textit{task\_ambiguous},\textit{stop\_bool},\textit{clarifying\_question}\}$.
Given both the image and the detected coordinates, the model learns to
cross-validate each region visually: when G-DINO points to a location
that does not contain the queried object, the model overrides the
detection and outputs \textsc{Stop}.

\paragraph{Training.}
We fine-tune Molmo-7B-D with LoRA~\cite{lora}
($r=8,\,\alpha=16$) for three epochs at learning rate $10^{-4}$ on
\datasettrain ($756$ augmented samples: horizontal flip,
brightness $\times 0.7$/$\times 1.3$, and coordinate perturbation with
$\sigma=5$\,px). To prevent the model from collapsing onto a
count-based shortcut, we include $64$ \emph{counterfactual pairs} that
share the same image and checkpoint coordinates but differ in $G_0$:
$G_0=\mathbf{x}_t$ labelled \textsc{Clear}/\textsc{Continue}, and
$G_0=\mathbf{x}_t+200$\,px labelled \textsc{Ambiguous}/\textsc{Ask}. We
additionally include $64$ ``no-$G_0$'' samples on $n=1$ trials labelled
by a count rule, which discourages a spurious
``moved$\,\Rightarrow\,$\textsc{Ask}'' shortcut when $G_0$ is
unavailable.

\subsection{Ablation design}

We isolate the contribution of each component using four systems:
\emph{AmbResVLM} (native two-step protocol; no $G_0$, no G-DINO, no
fine-tuning); \emph{Rule}$+G_0$ (count and distance rules over G-DINO
detections with $G_0$); \emph{FT-A} (fine-tuned with G-DINO but without
$G_0$); and \emph{FT-B} (full system). The AmbResVLM$\,\to\,$Rule$+G_0$
comparison isolates G-DINO (RQ1); FT-A$\,\to\,$FT-B isolates $G_0$
(RQ2); and Rule$+G_0\to$FT-B isolates fine-tuning (RQ3).

\section{Experiments}
\label{sec:exp}

\subsection{Setup}

\paragraph{Datasets.}
We train Stage~2 on \datasettrain ($70$ recorded trials
over the seven scenarios $S_1$--$S_7$). The training split contains
$756$ augmented samples; validation is held out at trial granularity
($84$ base samples) to prevent augmentation leakage. We evaluate on
\datasettest, an independent $70$-sample benchmark ($10$
trials per scenario) recorded in a separate session. Object compositions
span cup and cube targets together with red-box and box destinations;
each scenario introduces one class of mid-execution change at either
$C_1$ (pre-pick) or $C_2$ (pre-place), as summarized in
Table~\ref{tab:scenarios}.

\begin{table}[h]
\centering
\small
\caption{Scenarios in \datasettest (10 trials each).}
\label{tab:scenarios}
\begin{tabular}{clll}
\toprule
ID & Change & Gold decision & Checkpoint \\
\midrule
$S_1$ & no change                  & \textsc{Continue} & $C_1$ \\
$S_2$ & target count $+1$          & \textsc{Ask}      & $C_1$ \\
$S_3$ & target removed             & \textsc{Stop}     & $C_1$ \\
$S_4$ & target moved               & \textsc{Ask}      & $C_1$ \\
$S_5$ & destination count $+1$     & \textsc{Ask}      & $C_2$ \\
$S_6$ & destination removed        & \textsc{Stop}     & $C_2$ \\
$S_7$ & destination moved          & \textsc{Ask}      & $C_2$ \\
\bottomrule
\end{tabular}
\end{table}

\paragraph{Metrics.}
\emph{Decision accuracy} is the fraction of correct
\textsc{Continue}/\textsc{Ask}/\textsc{Stop} predictions. \emph{Miss
rate} is the fraction of non-\textsc{Continue} gold labels predicted as
\textsc{Continue} (an unsafe continuation). \emph{False-alarm rate}
(FAR) is the fraction of \textsc{Continue} gold labels predicted as
\textsc{Ask} or \textsc{Stop} (an unnecessary interruption). We report
miss rate and FAR separately because their associated costs are
asymmetric. \emph{Strict accuracy} additionally requires, on
\textsc{Ask} cases, that the generated question (a)~names the correct
object label and (b)~correctly characterizes the change type
(count increase: ``which $X$?''; movement: ``$X$ has moved''; absence:
``$X$ cannot be found''). \textsc{Continue} and \textsc{Stop} predictions
are scored on decision correctness alone.

\subsection{Main ablation}

Table~\ref{tab:main} reports headline metrics on
\datasettest.

\begin{table}[h]
\centering
\small
\caption{Main ablation on \datasettest ($70$ samples).
Strict accuracy is reported for fine-tuned variants that generate
clarifying questions.}
\label{tab:main}
\begin{tabular}{lccccccc}
\toprule
System & $G_0$ & G-DINO & FT & Dec.\,(\%) & Strict\,(\%) & Miss\,(\%) & FAR\,(\%) \\
\midrule
AmbResVLM            & N & N & N & 14.3 & --- & 100.0 & 0.0 \\
Rule\,$+G_0$         & Y & Y & N & 72.9 & --- & 38.3  & 0.0 \\
FT-A (no $G_0$)      & N & Y & Y & 68.6 & 61.4 & 36.7 & 0.0 \\
\textbf{FT-B (full)} & Y & Y & Y & \textbf{100.0} & \textbf{90.0} & \textbf{0.0} & \textbf{0.0} \\
\bottomrule
\end{tabular}
\end{table}

\paragraph{RQ1: G-DINO enables mid-execution detection.}
AmbResVLM attains $14.3\%$ decision accuracy ($100\%$ miss),
succeeding only on $S_1$ ($10/10$) and failing elsewhere because it
lacks count tracking, disappearance detection, and displacement
measurement. Augmenting the pipeline with G-DINO (Rule\,$+G_0$) raises
accuracy to $72.9\%$. The gain is most pronounced on count-based
scenarios: $S_2$ count increase (from $0/10$ to $10/10$),
$S_6$ destination absent (from $0/10$ to $6/10$), and the displacement
scenarios $S_4$/$S_7$ (from $0/10$ to $10/10$ each, enabled by $G_0$;
see RQ2).

\paragraph{RQ2: $G_0$ enables displacement detection.}
FT-A (fine-tuned without $G_0$) attains $68.6\%$ overall but $0/10$ on
both $S_4$ and $S_7$, because $n=1$ at a single position is
indistinguishable from \textsc{Clear} without a reference. FT-B reaches
$10/10$ on both, yielding a $+31.4$\,pp gain in overall accuracy.
Comparing Rule\,$+G_0$ ($72.9\%$, $S_4$/$S_7$ $10/10$) with FT-A
($68.6\%$, $S_4$/$S_7$ $0/10$) further shows that $G_0$ supplies
displacement information independently of model capacity.

\paragraph{RQ3: Fine-tuning corrects force-matching.}
Rule\,$+G_0$ fails entirely on \textsc{InvalidTarget} ($S_3$: $0/10$).
When the cup is removed, G-DINO force-matches to a visually similar
object at a displaced position, and the rule-based pipeline reads
$n=1$ at high distance as a moved target and outputs \textsc{Ask}.
FT-B reaches $10/10$ on $S_3$ by visually inspecting the G-DINO-reported
region and outputting \textsc{Stop} when no cup is present. The
$+27.1$\,pp improvement over Rule\,$+G_0$ quantifies the contribution
of cross-modal visual verification. Notably, FT-A also achieves
$10/10$ on $S_3$ without $G_0$, indicating that the correction is
driven by visual grounding rather than positional comparison.

\paragraph{Counterfactual diagnostic.}
To verify that the model relies on $G_0$ rather than a count-based
shortcut, we keep the image and the checkpoint detection coordinates
fixed while toggling $G_0$ between the current position (label
\textsc{Continue}) and a position shifted by $200$\,px (label
\textsc{Ask}). FT-B produces the appropriate decision for all $41$
single-detection trials, including every $S_4$ and $S_7$ trial, yielding
a memory reliance rate of $100\%$ ($41/41$) and confirming that the
decision depends on the stored $G_0$.

\section{Analysis}
\label{sec:analysis}

\begin{figure}[t]
\centering
\begin{subfigure}[t]{0.48\linewidth}
\centering
\includegraphics[width=\linewidth]{ambiguous_target_example.png}
\caption{$S_2$ (\textsc{AmbiguousTarget}). A second cup appears between
$t_0$ and $C_1$; FT-B outputs \textsc{Ask} with the question
``which cup do you mean?''.}
\label{fig:qual-amb}
\end{subfigure}
\hfill
\begin{subfigure}[t]{0.48\linewidth}
\centering
\includegraphics[width=\linewidth]{invalid_target_example.png}
\caption{$S_3$ (\textsc{InvalidTarget}). The cube is removed between
$t_0$ and $C_1$; FT-B outputs \textsc{Stop} with ``the cube cannot be
found'', despite G-DINO force-matching to a visually similar region.}
\label{fig:qual-inv}
\end{subfigure}
\caption{Representative FT-B outputs on \datasettest.
Each panel shows the initial scene ($t_0$, left) and the checkpoint
($C_1$, right). The decision and clarifying question are produced
jointly from the checkpoint image, the stored $G_0$, and the G-DINO
detections.}
\label{fig:qualitative}
\end{figure}

The $10$\,pp gap between decision accuracy and strict accuracy is
concentrated in $S_5$ (destination count increase); the remaining six
scenarios attain $100\%$ strict accuracy.

\paragraph{$S_5$: detector inconsistency degrades question quality.}
FT-B outputs \textsc{Ask} correctly on all $10$ trials but
mischaracterizes the underlying cause in $7$. The three trials on which
G-DINO successfully reports $n=2$ yield the correct question ``which
box should I place it on?''. On the remaining trials G-DINO reports a
single detection at a position different from $G_0$, leading the model
to infer either movement or absence. The count-increase signal is
suppressed at the detector stage, so strict accuracy on count-based
ambiguity is upper-bounded by detector consistency rather than by model
capacity.

\paragraph{Force-match correction example ($S_3$).}
On a representative trial, $G_0$ records the cup at $(268,147)$; after
the cup is removed, G-DINO force-matches the query ``cup'' to an
unrelated region at $(358,46)$, located $135$\,px from $G_0$.
Rule\,$+G_0$ reads $n=1$ with $\text{dist}>\tau$ and emits \textsc{Ask}.
FT-B inspects the reported region in the image, finds no cup, and
emits \textsc{Stop} with the explanation ``the cup cannot be found in
the scene''. The same behavior holds across all $10$ $S_3$ trials.

\paragraph{Label competition.}
When the target and destination categories share visual similarity
(e.g., ``cube'' and ``red box''), G-DINO may assign both detections to
the wrong label under a joint query. In $5$ of $10$ $S_3$ trials with
cube targets, the red box is misclassified as a cube and a spurious
$n=1$ detection reaches the evaluator. The fine-tuned model partially
recovers via visual verification, but the detector-level failure caps
strict accuracy on the affected trials. This failure mode originates
in the detector rather than in the evaluator.

\section{Conclusion}
\label{sec:conclusion}

We introduced \textbf{MemoryGrounding}, a two-stage framework that
addresses a gap in existing ambiguity-resolution methods: they evaluate
the scene only at task initiation and cannot respond to dynamic changes
during execution. Controlled ablation demonstrates that (i)~existing
VLM-based detection ignores count and position
($14.3\%$ for AmbResVLM), and integrating G-DINO restores quantitative
scene perception; (ii)~storing the initial grounding $G_0$ provides
the temporal memory required for displacement detection ($+31.4$\,pp
over FT-A); and (iii)~fine-tuning a Molmo-7B evaluator enables
cross-modal visual verification that overrides spurious open-vocabulary
detections, whereas rule-based systems fail on every target-removal
trial. The full system attains $100\%$ decision accuracy with $0\%$
miss rate and FAR, $90.0\%$ strict accuracy on
\datasettest, and a memory reliance rate of $100\%$ under
a counterfactual diagnostic that varies $G_0$ at fixed image and
detections.

\paragraph{Limitations.}
The $10$\,pp decision-strict gap is concentrated in $S_5$ and is driven
by detector inconsistencies that suppress the count-increase signal.
Label competition between visually similar categories under joint
queries further caps performance. The benchmark
\datasettest is saturated by FT-B, which limits our
ability to differentiate richer multi-checkpoint episode memory on this
dataset. The Stage~1 thresholds are tuned empirically and may require
re-calibration across camera setups, object scales, and robot speeds.
Our analysis also assumes that $G_0$ is correctly established at $t_0$;
integrating an initial-scene disambiguation front-end is a natural
extension.

\paragraph{Future work.}
We highlight four directions. First, multi-checkpoint episode memory
generalizing $G_0$ to a full execution history. Second, end-to-end
evaluation in which Stage~1 dynamically supplies trigger points on
continuous video. Third, detector robustness via re-ranking or visual
re-identification to mitigate force-matching and label competition.
Fourth, comparison against strong general-purpose VLMs (e.g., GPT-4V,
Gemini, Claude) prompted with the same $G_0$-conditioned structure, and
generalization to larger object families, longer-horizon manipulation,
and more complex spatial relations.

\bibliographystyle{unsrt}
\bibliography{references}

\newpage
\appendix

\section{Appendix: Team contributions}

Seongmin Lee led system design and implemented the Stage~1 FSM trigger
and the rule-based pipeline (Rule$+G_0$); performed all data collection
and recording of \datasettrain and \datasettest; and contributed to
evaluation and writing.
Sanghyuk Ko designed and implemented the Stage~2 evaluator, including
LoRA fine-tuning of Molmo-7B-D, counterfactual augmentation strategy,
and the structured JSON decision policy; led ablation analysis and
paper writing.

\end{document}
